% statement-two-page.tex — layout stress for the statement continuation header,
% Page N of M folio, ordinary headings, and natural prose page breaking.
\documentclass[type=research]{careerdossier-statement}
\CDossierSetup{
  name={Ada Lovelace},
  email={ada.lovelace@example.com},
  phone={+1 416 555 0142},
  website={example.com/ada},
  affiliation={Example University},
  scholar={scholar.google.com/citations?user=ada-example},
  orcid={0000-0002-1825-0097}
}
\CDossierStatementSetup{
  running-title={Computational Reliability},
  subtitle={Transparent methods for trustworthy scientific computing},
  application-context={Application for Assistant Professor of Computational Science},
  application-id={APP-2026-0042}
}
\begin{document}
\MakeCDossierStatementHeader

\section*{Research Vision}

My research asks how computational results can remain understandable,
reproducible, and useful after the immediate project that produced them has
ended. I develop numerical methods and research-software practices that expose
their assumptions, preserve the path from evidence to conclusion, and allow
scientists from neighbouring fields to test and extend the work.

The central premise is that reliability is not a final verification step. It is
an organizing principle for model design, implementation, documentation, and
communication. A method is most valuable when its mathematical guarantees, its
software realization, and its limits can be inspected together.

\section*{Prior and Current Contributions}

My first line of work develops adaptive algorithms whose error estimates remain
meaningful when data quality and computational resources vary. I combine
analysis with executable experiments, using small reference implementations to
make each theoretical claim testable. This approach has helped collaborators
distinguish uncertainty arising from a physical model from error introduced by
discretization or implementation choices.

A second line of work concerns durable computational records. I design workflows
that connect source data, environment descriptions, tests, and figures without
requiring readers to reconstruct an undocumented sequence of manual steps. In
collaborative projects, these records have supported method comparison,
onboarding, and independent replication. They also provide students with a
concrete model of research accountability: a result is incomplete until another
person can understand how it was produced.

These projects share a practical commitment to accessibility across expertise.
I write software interfaces and technical explanations for readers who know the
scientific question but may not know the implementation language or numerical
method. This requires choosing meaningful defaults, naming uncertainty clearly,
and separating essential decisions from incidental technical detail.

\section*{Future Research Programme}

The next phase of my programme has three connected aims. First, I will develop
verification methods for computational pipelines that combine learned and
mechanistic models. The goal is to identify which guarantees can be preserved,
which assumptions must be measured empirically, and how uncertainty should be
reported when components change over time.

Second, I will study reproducibility as a property of research infrastructure
rather than an after-the-fact archive. This work will test lightweight ways to
record environments, data transformations, and validation decisions during
ordinary scientific practice. The expected outputs include open reference
implementations, reusable evaluation datasets, and guidance that small research
groups can adopt without a dedicated engineering team.

Third, I will build teaching and mentorship into the research programme.
Students will participate in paired analysis and implementation reviews, learn
to design tests from mathematical claims, and publish documentation for readers
outside the immediate project. This structure gives students ownership of
substantive questions while making the standards of reliable work explicit.

\section*{Institutional Fit and Trajectory}

Example University's strengths in computational science, applied mathematics,
and open research create a productive setting for this agenda. I would seek
collaborations with researchers whose domain questions stress current methods,
and I would contribute shared software and training that make those
collaborations easier to sustain. Over the next five years, I aim to establish a
research group known for connecting rigorous numerical reasoning with tools and
practices that other scientists can genuinely reuse.

\end{document}
